\documentclass[a4paper,12pt]{article}

\usepackage[utf8]{inputenc}
\usepackage[russian]{babel}
\usepackage{geometry}
\usepackage{graphicx}
\usepackage{hyperref}
\usepackage{listings}
\usepackage{xcolor}
\usepackage{float}
\usepackage{caption}

\usepackage[T2A]{fontenc}
\usepackage[utf8]{inputenc}
\usepackage[russian]{babel}
\usepackage{geometry}
\usepackage{listings}
\usepackage{xcolor}

\geometry{margin=2cm}

\lstset{
basicstyle=\ttfamily\small,
breaklines=true,
frame=single,
columns=fullflexible
}

% Оформление листингов
\lstset{
    basicstyle=\ttfamily\small,
    keywordstyle=\color{blue},
    commentstyle=\color{gray},
    stringstyle=\color{red},
    numbers=left,
    numberstyle=\tiny,
    breaklines=true,
    frame=single,
    tabsize=2
}

\begin{document}

\pagestyle{empty}

%-------------------- ТИТУЛЬНЫЙ ЛИСТ --------------------

\begin{center}
Санкт-Петербургский национальный исследовательский университет\\
информационных технологий, механики и оптики

\vfill
Факультет прикладной информатики

\vfill
\textbf{Web-программирование}

Лабораторная работа №4
\end{center}

\vfill

\begin{flushright}
Выполнил: Дармороз М.Д.\\
Проверил: Богатыренков С.С.
\end{flushright}

\vfill

\begin{center}
Санкт-Петербург\\
2026
\end{center}

\newpage
\pagestyle{plain}

\section{Введение}

Целью данной лабораторной работы является разработка веб-приложения на
Node.js, работающего по протоколу HTTPS, которое реализует два маршрута:
\begin{itemize}
  \item \texttt{/login} --- возвращает логин студента в системе;
  \item \texttt{/id/\{N\}} --- выполняет GET-запрос к внешнему API
        \url{https://nd.kodaktor.ru/users/\{N\}} без заголовка
        \texttt{Content-Type} и возвращает значение поля \texttt{login}
        из полученного ответа.
\end{itemize}

В качестве стека технологий выбраны платформа Node.js и фреймворк
Express.js. HTTP-запрос к внешнему ресурсу выполняется встроенным
модулем \texttt{https} для обеспечения полного контроля над
отправляемыми заголовками.

\section{Ход работы}

\subsection{Установка зависимостей}

Был инициализирован пустой проект и установлены необходимые пакеты:

\begin{lstlisting}[language=bash, caption={Инициализация проекта}]
npm init -y
npm install express
\end{lstlisting}

Пакет \texttt{express} используется для маршрутизации и обработки
HTTP-запросов. Модуль \texttt{https} является встроенным и не требует
установки.

\subsection{Структура проекта}

Файловая структура проекта выглядит следующим образом:

\begin{lstlisting}[language=bash, caption={Структура проекта}]
weblab4/
  app.js          % исходный код приложения
  package.json    % конфигурация проекта и зависимости
  node_modules/   % установленные пакеты
\end{lstlisting}

\subsection{Исходный код}

Основной файл приложения --- \texttt{app.js}. Его код представлен
ниже:

\begin{lstlisting}[language=JavaScript, caption={app.js --- основной сервер}, label={lst:app}]
const express = require('express');
const https = require('https');

const app = express();
const PORT = process.env.PORT || 3000;

app.get('/login', (req, res) => {
  res.send('amaenai1_');
});

app.get('/id/:n', (req, res) => {
  const { n } = req.params;
  const url = new URL(`https://nd.kodaktor.ru/users/${n}`);

  const options = {
    hostname: url.hostname,
    path: url.pathname,
    method: 'GET',
  };

  https.request(options, (response) => {
    let data = '';
    response.on('data', chunk => data += chunk);
    response.on('end', () => {
      try {
        const parsed = JSON.parse(data);
        res.send(parsed.login);
      } catch {
        res.status(500).send('Error parsing response');
      }
    });
  }).on('error', (err) => {
    res.status(500).send(err.message);
  }).end();
});

app.listen(PORT, () => {
  console.log(`Server running on http://localhost:${PORT}`);
});
\end{lstlisting}

\subsection{Описание работы программы}

При запуске сервер начинает прослушивать порт, указанный в
переменной окружения \texttt{PORT} (или порт 3000 по умолчанию).
Приложение обрабатывает два типа запросов:

\textbf{Маршрут \texttt{/login}.} При GET-запросе к этому адресу
сервер немедленно возвращает строку \texttt{amaenai1\_} --- логин
студента в системе Moodle. Реализация тривиальна: обработчик
вызывает \texttt{res.send()} с готовой строкой.

\textbf{Маршрут \texttt{/id/\{N\}}.} Данный маршрут принимает
параметр \texttt{N} из URL-адреса. Алгоритм его работы:

\begin{enumerate}
  \item Параметр \texttt{N} извлекается из пути с помощью
        \texttt{req.params}.
  \item Формируется целевой URL вида
        \texttt{https://nd.kodaktor.ru/users/N}.
  \item Создаётся объект \texttt{options}, содержащий только
        имя хоста и путь. Поле \texttt{headers} отсутствует, что
        гарантирует, что заголовок \texttt{Content-Type} не будет
        отправлен в запросе.
  \item Выполняется GET-запрос методом \texttt{https.request()}.
  \item Ответ собирается по частям через событие \texttt{data}.
  \item После получения всего тела ответа (событие \texttt{end})
        строка парсится в JSON. Если парсинг успешен, клиенту
        возвращается значение поля \texttt{login}. В противном
        случае возвращается ошибка 500.
  \item При возникновении ошибки сетевого уровня (событие
        \texttt{error}) клиенту также возвращается ошибка 500 с
        текстом сообщения.
\end{enumerate}

Ключевая особенность --- использование встроенного модуля
\texttt{https} вместо библиотек наподобие \texttt{axios}. Это
позволяет тонко управлять заголовками: при использовании
\texttt{axios} по умолчанию устанавливается заголовок
\texttt{Content-Type: application/json}, что противоречит
требованию задания. Модуль \texttt{https.request()}, напротив,
не добавляет никаких заголовков, если они не указаны явно.

\subsection{Деплой на Render}

Для публикации приложения в интернете с поддержкой HTTPS была
использована платформа Render. Процесс деплоя включает следующие
шаги:

\begin{enumerate}
  \item Код загружен в репозиторий на GitHub.
  \item В личном кабинете Render создан новый Web Service,
        подключённый к этому репозиторию.
  \item Указана команда запуска: \texttt{node app.js}.
  \item Выбран бесплатный тарифный план.
\end{enumerate}

После сборки и запуска приложение становится доступно по адресу
\texttt{https://weblab4.onrender.com}. Платформа Render
автоматически выдаёт SSL-сертификат и проксирует все запросы по
HTTPS, поэтому дополнительная настройка не требуется.

\section{Заключение}

В ходе лабораторной работы было разработано веб-приложение на
Node.js с использованием фреймворка Express.js. Приложение
реализует два маршрута: \texttt{/login}, возвращающий логин
студента, и \texttt{/id/\{N\}}, выполняющий прокси-запрос к
внешнему API и возвращающий поле \texttt{login} ответа.

Удалось выполнить все требования задания:
\begin{itemize}
  \item приложение работает по протоколу HTTPS (за счёт
        платформы Render);
  \item заголовок \texttt{Content-Type} не отправляется в
        запросе к \texttt{nd.kodaktor.ru};
  \item код корректен и протестирован.
\end{itemize}

В качестве дополнительного результата получен опыт деплоя
Node.js-приложений на платформу Render, включая настройку
автоматического HTTPS.

\end{document}
